\section{Detection on ordinary spectra}
\label{sec:detection}

We now turn the coefficient relation into a statement about the natural
localization map.  Keep $G=P_3\times P_2$ and
$H=P_2\times\overline Z$, and put
\[
 \Delta=\Gal(k/\F_p),\qquad
 G_3=P_3\rtimes\Delta,\qquad
 \mathcal D=\bigl((\widehat B^{\square})^{hG}\bigr)(*).
\]
Thus $\Delta$ has order six.  The symbol $\mathcal D$ denotes an ordinary spectrum:
the continuous fixed points are formed in solid spectra and then evaluated
at the point.

\begin{lemma}[The ordinary local detector]\label{lem:local-detector}
There is a natural map of ordinary spectra
\begin{equation}\label{eq:ordinary-detector}
 g:L_{K(3)}S\longrightarrow\mathcal D.
\end{equation}
The spectrum $\mathcal D$ is $K(2)$-local, so $g$ factors through the canonical
localization unit:
\begin{equation}\label{eq:detector-factorization}
 L_{K(3)}S\xrightarrow{\eta}L_{K(2)}L_{K(3)}S
       \xrightarrow{\overline g}\mathcal D.
\end{equation}
\end{lemma}

\begin{proof}
Use the ordinary descent identification in
\cref{prop:completed-theory} to form the composite
\begin{equation}\label{eq:detector-construction}
\begin{split}
 L_{K(3)}S
 &\simeq\bigl((E_3^{\square})^{hG_3}\bigr)(*)
 \longrightarrow\bigl((E_3^{\square})^{hP_3}\bigr)(*)\\
 &\longrightarrow\bigl((E_3^{\square})^{hG}\bigr)(*)
 \longrightarrow\bigl((\widehat B^{\square})^{hG}\bigr)(*).
\end{split}
\end{equation}
The first arrow restricts the group action, the second is inflation along
$G\to P_3$, and the last is induced by the equivariant map
\eqref{eq:completed-map}.  In the middle term with group $G$, the action on
$E_3^{\square}$ is through $P_3$. Restriction and inflation are the
natural maps of continuous fixed points. The source equivalence and its ordinary unit are the point-valued
form of \cite[Proposition~4.6.9]{BMR2026} proved in
\cref{prop:completed-theory}; the last map uses the local split-section
completed map of \cref{lem:completed-oriented-map} and that proposition.

For every profinite set $T$, the evaluation
$\widehat B^{\square}(T)$ is the ordinary height-two Lubin--Tate theory
described in \cref{prop:completed-theory}, with the Noetherian witness
verified there; hence it is $K(2)$-local
\cite[Proposition~4.4.2]{BMR2026}.
The bar construction gives
\[
 \mathcal D\simeq\operatorname{Tot}\bigl(\widehat B^{\square}(G^{\bullet})\bigr).
\]
Ordinary local spectra are closed under limits, so $\mathcal D$ is
$K(2)$-local.  This also follows from the
pointwise criterion for solid locality
\cite[Definition~4.5.7 and the following paragraph]{BMR2026}.
The universal property of ordinary $K(2)$-localization now gives
\eqref{eq:detector-factorization}.
\end{proof}

\subsection{Convergence before rationalization}

We use integral continuous cohomology of the solid homotopy modules in
\cref{lem:solid-resolutions}.  For such a module $M$, recall that
$H^s(A,M)$ denotes the point evaluation of $R^s\Gamma(A,M)$.
In particular, $[1/p]$ below is applied to this integral group after it
has been computed.  For $M=\Z_p$, the natural comparison
\eqref{eq:group-cochain-bridge} identifies these groups with the
finite continuous-cochain derived-limit convention of
\eqref{eq:integral-group-cochains}.

The compact groups needed here have finite cohomological dimension in
this coefficient category.  Indeed, if the degree-$n$ division algebra
defining $P_n$ contained an element of order $p$, it would contain the
field $\Q_p(\zeta_p)$: the nontrivial element has the irreducible
cyclotomic minimal polynomial of degree $p-1$.  Its reduced
characteristic polynomial has degree $n$ and annihilates it.
Since $p-1>n$ for $n=2,3$ and $p\geq5$, this is impossible.
Thus neither $P_2$ nor $P_3$ has $p$-torsion.  The same holds for
$G$ and $H$.  An order-$p$ element of $G_3$ would project trivially to
$\Delta$, since $p\nmid6$, and hence would lie in $P_3$.  Thus
\cref{lem:solid-resolutions} gives a uniform finite cohomological bound
for each of these groups.  We only need the existence of a common bound
$d$, independent of the homotopy degree of the coefficient module.

\begin{lemma}[Finite-amplitude descent]\label{lem:descent-convergence}
For the continuous actions used in \eqref{eq:detector-construction},
there are natural, strongly convergent spectral sequences
\begin{equation}\label{eq:integral-descent-ss}
 E_2^{s,t}=H^s(A,\pi_tX)
       \ \Longrightarrow\ \pi_{t-s}\bigl((X^{hA})(*)\bigr),
 \qquad 0\leq s\leq d.
\end{equation}
They give, after exact rationalization, natural strongly convergent
spectral sequences with $E_2$-terms $H^s(A,\pi_tX)[1/p]$ and abutments
$\pi_{t-s}((X^{hA})(*))[1/p]$.
This statement applies to the unbounded even-periodic theories here.
\end{lemma}

\begin{proof}
Let $F(X)=(X^{hA})(*)$.  This is an exact functor preserving limits and
upper homotopy bounds.  For the latter assertion, coconnective objects
remain coconnective under evaluation and limits.  Although $X$ need not
be a module spectrum over $H\Z_p$, each Postnikov layer $HM$ is one.
The forgetful functor from module spectra preserves limits, so its
continuous fixed points agree with the derived module invariants.
On this Eilenberg--Mac Lane object, the finite
resolution in \cref{lem:solid-resolutions} gives
\[
 \pi_{-s}F(HM)=H^s(A,M),\qquad
 F(HM)\in\mathrm{Sp}_{[-d,0]}.
\]
The solid spectra are used with the hypersheaf convention of
\cite[Section~4.5, especially Proposition~4.5.5]{BMR2026}.
Truncations remain solid, since they retain or kill each solid homotopy
module.  Evaluation on an extremally disconnected profinite set is
$t$-exact: homotopy sheafification evaluated on this projective object
agrees with pointwise homotopy.  It therefore commutes with truncation.
These evaluations detect equivalences, so Postnikov completeness of
ordinary spectra gives
$X\simeq\varprojlim_n\tau_{\leq n}X$.

Put $T_n=F(\tau_{\leq n}X)$.  Exactness identifies the successive fiber as
\[
 \operatorname{fib}(T_{n+1}\longrightarrow T_n)
   \simeq F\bigl(\Sigma^{n+1}H\pi_{n+1}X\bigr)
   \in\mathrm{Sp}_{[n+1-d,n+1]}.
\]
Consequently $\pi_qT_{n+1}\to\pi_qT_n$ is an isomorphism for
$n\geq q+d$; the analogous tower in degree $q+1$ is also eventually
constant.  The Milnor exact sequence therefore has zero derived-limit
term, and gives
\begin{equation}\label{eq:finite-postnikov-window}
 \pi_qF(X)\cong\pi_qF(\tau_{\leq q+d}X).
\end{equation}
The infinitely negative tail causes no further issue: since
$F(\tau_{\leq q-1}X)$ is concentrated in degrees at most $q-1$, the
fiber sequence for this lower cutoff gives
\[
 \pi_qF(\tau_{[q,q+d]}X)
     \xrightarrow{\ \sim\ }\pi_qF(\tau_{\leq q+d}X).
\]
Hence each abutment degree is computed by a finite Postnikov interval.

The exact couple of this Postnikov construction has the $E_2$-term in
\eqref{eq:integral-descent-ss}.  Its filtration in degree $q$ is
\[
 \operatorname{Fil}^{s}\pi_qF(X)
 =\ker\!\left(\pi_qF(X)\longrightarrow
              \pi_qF(\tau_{\leq q+s-1}X)\right),
 \quad 0\leq s\leq d+1.
\]
It starts with the whole group, ends in zero by
\eqref{eq:finite-postnikov-window}, and has associated graded
$E_\infty^{s,q+s}$.  This proves the asserted strong convergence and its
naturality.  Tensoring the exact couple and these finite filtrations with
$\Q$ is exact.  It therefore gives the rational spectral sequence and
its stated abutment; no exchange of rationalization with an infinite
totalization is involved.  The groups here are $\Z_p$-modules, so this
tensor operation agrees with $[1/p]$.
No uniform torsion exponent as $t$ varies is required: each fixed
homotopy degree has the finite filtration just constructed.
\end{proof}

\subsection{The surviving coefficient row}

\begin{lemma}[Central weights]\label{lem:central-weights}
For $t\ne0$ and every $s$, one has
\[
 H^s(G_3,\pi_tE_3^{\square})[1/p]=0,
 \qquad
 H^s(G,\pi_t\widehat B^{\square})[1/p]=0.
\]
\end{lemma}

\begin{proof}
Odd homotopy modules vanish.  For $t=2m$, scalar automorphisms
$a\in1+p\Z_p$ lift canonically to scalar automorphisms of the universal
deformation.  They fix its deformation ring, and act by $a^m$ on its
$m$th invariant-differential line $\pi_{2m}E_n^{\square}$.
This is the central character used in
\cite[Lemma~2.6.1]{BSSW2025}; replacing the periodicity convention by its
dual replaces $m$ by $-m$ and has the same consequence.

Choose $b=1+p$.  It is central in $G_3$, since the finite Galois action
fixes these scalars.  A central group element acting on a coefficient
module induces the identity on its group cohomology.  This follows
already from the bar resolution: its coefficient action is chain
homotopic to the action by conjugation, which is the identity for a
central element.  Consequently $(b^m-1)$ annihilates
$H^s(G_3,\pi_{2m}E_3^{\square})$.  For $m\ne0$ this is a nonzero
$p$-adic scalar, and is invertible after $p$ is inverted.  This proves
the first vanishing.

For the target, use the natural coefficient descent
\eqref{eq:coefficient-descent} to identify its cohomology with
$H^s(H,\pi_tE_2^{\square})$.  The factor $\overline Z$ acts trivially
on these coefficients, and the central scalars in $P_2$ act by the same
character $b^m$.  The preceding argument applies to $H$, and proves
the second vanishing.
\end{proof}

\begin{proposition}\label{prop:nonzero-detector}
The ordinary map $g$ of \eqref{eq:ordinary-detector} is nonzero on
rational homotopy in degree $-3$:
\[
 \pi_{-3}L_0g:\pi_{-3}L_0L_{K(3)}S
                    \longrightarrow\pi_{-3}L_0\mathcal D
 \quad\text{is nonzero}.
\]
\end{proposition}

\begin{proof}
First extend the source primitive to the extended stabilizer.
By \cref{prop:group-classes}, $e_{P_3}$ is represented under rational
Lie comparison by a nonzero multiple of
\[
 (X,Y,Z)\longmapsto\tr_{\mathrm{red}}\bigl(X[Y,Z]\bigr).
\]
In the standard presentation of the degree-three division algebra,
take its maximal unramified subfield $\Q_{p^3}$ and an element $\Pi$
such that
\[
 \Pi x=\sigma(x)\Pi\quad(x\in\Q_{p^3}),\qquad \Pi^3=p,
\]
where $\sigma$ is arithmetic Frobenius of the unramified cubic
extension $\Q_{p^3}/\Q_p$, inducing $z\mapsto z^p$ on the residue
field. Its extension fixing $\Pi$ equals conjugation
$y\mapsto\Pi y\Pi^{-1}$, as is checked on the generators
$\Q_{p^3}$ and $\Pi$. It therefore preserves the reduced-trace form.
Its action on $P_3$ factors through the order-three quotient of
$\Delta$; the order-two kernel also fixes the $\Z_p$-coefficients.
Naturality of Lie comparison therefore makes $e_{P_3}$ invariant under
$\Delta$.  Since $6$ is a unit in $\Z_p$, finite-group averaging and
extension descent give
\[
 H^3(G_3,\Z_p)[1/p]
    \xrightarrow{\ \sim\ }
       \bigl(H^3(P_3,\Z_p)[1/p]\bigr)^\Delta.
\]
Let $\widetilde e_3$ be the class restricting to $e_{P_3}$, and let
$z_3$ be its image under the unit-coefficient map
$\Z_p\to\pi_0E_3^{\square}$.

Under \eqref{eq:group-cochain-bridge}, restriction from $G_3$ to $P_3$
and inflation along $G\to P_3$ on the constant classes are the
derived limits of their finite continuous-cochain pullbacks.  These
are the group maps used in \eqref{eq:constant-cochain-naturality}.
All coefficient maps in the following square are the maps induced by
\eqref{eq:detector-construction}:
\[
\begin{tikzcd}[column sep=large]
 H^3(G_3,\Z_p)[1/p]
   \arrow[r] \arrow[d,"{\mathrm{res},\,\mathrm{inf}}"']
 & H^3(G_3,\pi_0E_3^{\square})[1/p]\arrow[d]\\
 H^3(G,\Z_p)[1/p]\arrow[r]
 & H^3(G,\pi_0\widehat B^{\square})[1/p].
\end{tikzcd}
\]
On bar cochains, the right vertical map is restriction, pullback along
$G\to P_3$, and the coefficient map of
$E_3^{\square}\to\widehat B^{\square}$, in that order.
Its composite with the source coefficient unit is the target unit,
which is also $\Z_p\to W(\widehat L)\to\pi_0\widehat B^{\square}$.
Here the first map is the integral limit of the finite Witt units in
\eqref{eq:finite-witt-constant-square}; the map from full Witt vectors
to deformation coefficients is applied only after that limit.
The source descent identification used here is the canonical one in
\cref{prop:completed-theory}.
In particular the image of $z_3$ is exactly the class called $c_3$ in
\cref{prop:coefficient-transport}: it is obtained from $e_{P_3}$ by
projection and the actual constant-coefficient inclusion.  That
proposition and \cref{lem:constant-nonvanishing} give
\[
 c_3=a\,c_2\ne0,\qquad a\in\Q_p^{\times}.
\]
Thus the map between the rational descent $E_2$-terms is nonzero in
bidegree $(s,t)=(3,0)$.

Apply \cref{lem:descent-convergence} to the source and target of $g$.
By \cref{lem:central-weights}, only $t=0$ survives after rationalization.
The differential has bidegree $(r,r-1)$, so every rational differential
vanishes.  For total homotopy degree $-3$, the finite filtration has
just one possibly nonzero graded piece, namely $(s,t)=(3,0)$.
The nonzero map on this piece is therefore a nonzero map on the
abutments.  These abutments are the ordinary rational homotopy groups
asserted in the proposition.
\end{proof}

\begin{proof}[Proof of \cref{thm:main}]
Rationalize \eqref{eq:detector-factorization}.  If the canonical map
$L_0L_{K(3)}S\to L_0L_{K(2)}L_{K(3)}S$ were zero on $\pi_{-3}$,
its composite $L_0g$ would also be zero there.  This contradicts
\cref{prop:nonzero-detector}.  Every step above applies to each prime
$p\geq5$, proving the theorem.
\end{proof}
